This paper presents an empirical study on integrating conversational AI (specifically ChatGPT) into requirements engineering (RE) education within an undergraduate software engineering course. The authors designed an experiment involving 20 student groups who developed software projects and used ChatGPT during the requirements specification phase under two different processes: one where students first created requirements manually before refining them with ChatGPT (Process A), and another where students directly used ChatGPT from the start (Process B). The study evaluates the quality of requirements artifacts, prompting strategies, and student effort. Results show that both approaches yield comparable outcomes, with average artifact quality around 80\%, but achieving higher-quality results requires significant effort in prompt design, critical evaluation, and refinement. The study also highlights a strong correlation between prompt quality and output quality, emphasizing the importance of structured, detailed, and context-rich prompts.

Beyond performance, the paper explores how students interact with conversational AI and their perceptions of its usefulness. Students primarily used ChatGPT for idea generation, refinement, and rewriting, with “concretization” (refining existing ideas) being the most effective use case. However, the study identifies several challenges: students often lacked prompting proficiency, were sometimes overly influenced by ChatGPT's confident responses, and frequently accepted generic or incorrect outputs without sufficient critical analysis. The findings suggest that while conversational AI can enhance productivity and support learning, its effective use requires training in prompt engineering and critical thinking. Overall, the paper argues for a shift in software engineering education toward teaching students how to collaborate with AI tools effectively rather than treating them as simple automation aids.

\section{Motivation}
The introduction motivates the study by highlighting the rapid and disruptive emergence of conversational AI tools such as ChatGPT and GitHub Copilot, which have significantly reshaped computing education and software engineering practices. These tools introduce both opportunities and challenges for instructors and students, requiring a rethinking of traditional pedagogical approaches \cite{openai2023chatgpt,githubcopilot2021,denny2023prompt,becker2023ai,prather2023ai,kasneci2023chatgpt}. The authors emphasize that resisting this shift is impractical, echoing prior arguments that adaptation to AI-assisted workflows is inevitable \cite{kasneci2023chatgpt}. Consequently, there is a pressing need to understand how such tools can be systematically integrated into education while preserving learning objectives and fostering critical thinking.

Another key motivation stems from the gap between the availability of powerful AI tools and the lack of structured guidance on how to use them effectively in software engineering tasks, particularly in requirements engineering. While prior work has explored AI-assisted coding and educational applications, fewer studies examine how conversational AI can support early-stage development activities such as requirements elicitation and specification. This is especially important because requirements engineering is a cognitively demanding and collaborative process that benefits from creativity, precision, and iterative refinement. The authors thus position their work as an effort to explore how AI can augment, rather than replace, human reasoning in these tasks, aligning with broader calls to adapt teaching practices to the AI era \cite{denny2023prompt,becker2023ai}.

Finally, the study is motivated by the need to evaluate not only the effectiveness of AI tools but also their impact on student behavior, learning, and decision-making. Prior research suggests that AI tools can influence how students approach problem-solving, sometimes leading to over-reliance or reduced critical engagement \cite{prather2023ai,kasneci2023chatgpt}. Therefore, the authors aim to investigate how students interact with conversational AI in a controlled educational setting, what strategies they adopt, and how these interactions affect the quality of their outputs. This motivation reflects a broader educational goal: preparing students to collaborate effectively with AI systems by developing skills in prompt engineering, evaluation, and critical thinking.

\section{Methodology}
\subsection{Course Setup}
The study was conducted in the context of an upper-level undergraduate software engineering course where students worked in 20 groups of four to develop large-scale projects involving an Android client and a Node.js backend. The course followed an iterative development process, requiring multiple deliverables such as requirements, design, implementation, testing, and documentation. A key pedagogical objective was to systematically integrate conversational AI while fostering critical thinking, allowing students to use ChatGPT in a controlled manner rather than as an unrestricted tool.

To evaluate AI-assisted requirements engineering, the instructors designed two structured workflows: Process A, where students first created requirements manually and then refined them using ChatGPT, and Process B, where students used ChatGPT from the beginning. Students were required to submit final artifacts, intermediate drafts, and complete ChatGPT interaction logs, enabling detailed analysis of both outcomes and processes. This setup ensured that the study could compare human-first vs. AI-first workflows while capturing how students iteratively engaged with the tool.

\subsection{RQ1: Effectiveness of ChatGPT Interaction Processes}
This research question investigates how different interaction modes (Process A vs. Process B) affect the quality of requirements artifacts. The methodology involved evaluating student submissions using predefined grading criteria to ensure consistency across groups. The study compares outputs generated through manual-first refinement versus direct AI-driven generation, enabling a structured assessment of whether early human reasoning leads to better outcomes than immediate AI reliance.

The evaluation focuses on determining whether either process leads to significantly better results or whether both approaches converge to similar quality levels. A key aspect of this analysis is understanding the trade-off between human effort and AI assistance, particularly whether AI-first approaches reduce effort or simply shift it toward prompt engineering and refinement. The study thus examines not only final outputs but also the process efficiency and effectiveness of each interaction mode.

\subsection{RQ2: Prompting Strategies and Their Impact}
This research question explores how students construct prompts and how these strategies influence the quality of ChatGPT-generated outputs. The methodology includes analyzing the full conversation logs submitted by students to identify patterns in prompting behavior, such as level of detail, structure, ambiguity, and use of examples. The authors particularly examine how different prompting styles correlate with improvements or deficiencies in the resulting requirements.

A central finding investigated here is the relationship between prompt quality and output quality, with emphasis on how vague prompts (e.g., “improve this”) lead to misaligned or generic responses, while structured prompts produce more relevant results. The study also evaluates students' prompting proficiency, identifying common shortcomings such as lack of context, insufficient constraints, and unclear expectations. This analysis highlights the importance of prompt engineering as a learned skill in AI-assisted development workflows.

\subsection{RQ3: Student Effort and Interaction Behavior}
This research question examines how students interact with ChatGPT and the effort required to achieve high-quality results. The methodology involves analyzing both qualitative and quantitative data from interaction logs to understand how students allocate time across activities such as brainstorming, prompting, and validating outputs. The study pays particular attention to how effort differs across interaction modes and how it influences final artifact quality.

A key aspect of this analysis is identifying behavioral patterns, including students being overly influenced by ChatGPT's confident responses and sometimes accepting incorrect suggestions without sufficient critique. The study also evaluates how much effort is needed to move beyond mid-range quality outputs (~80\%), emphasizing that higher-quality results require significant investment in refinement and critical assessment. This highlights the role of human oversight and critical thinking in effective human-AI collaboration.

\section{Key Findings}

\subsection{RQ1: Process A vs. B}
The results show that students achieved comparable quality across all processes, with average scores of 82\% ($A^+$), 75\% ($A^-$), and 77\% (B), and no statistically significant difference between them . Interestingly, half of the top-performing projects came from each process, indicating that neither human-first nor AI-first workflows dominate in terms of final outcomes. However, a key behavioral difference emerges: students in Process A (especially $A^+$) were more likely to edit ChatGPT outputs, suggesting stronger engagement and critical refinement compared to Process B.

Another important observation is that many students in Process A (7/20 groups) did not actually provide their initial human-written specifications to ChatGPT, instead using it only for idea generation. This effectively made their workflow closer to Process B, highlighting inconsistencies in how students applied the intended methodology. Overall, while both processes can yield similar results, Process A encourages more active editing and critical interaction, whereas Process B tends to involve more direct reliance on AI-generated content.

\subsection{RQ2: Effort}
The study finds that achieving high-quality outputs with ChatGPT requires substantial effort beyond initial generation, particularly in prompt design, iterative refinement, and validation of outputs . While students can relatively easily reach a mid-level quality (~80\%), improving beyond this threshold demands significant investment in brainstorming, crafting detailed prompts, and critically assessing AI responses. This highlights that AI does not eliminate effort but rather shifts it toward more cognitive and interaction-based tasks.

Additionally, the analysis shows that prompt quality is strongly tied to the effort invested, with better prompts requiring more time, more detailed instructions, and clearer structure. The expert investigator experiment reinforces this finding, demonstrating that high-quality outputs can be consistently achieved with well-structured prompts, but only when sufficient effort is spent on designing them. Thus, the key takeaway is that effective use of ChatGPT is effort-intensive and skill-dependent, rather than a shortcut to high-quality results.

\subsection{RQ3: Usage Patterns}
Students primarily used ChatGPT for ideation, refinement, and formalization of requirements, with refinement and rewriting being the most commonly adopted uses . This suggests that ChatGPT is most effective when students already have partially developed ideas, allowing the tool to enhance clarity and structure rather than generate content from scratch. Different students exhibited different preferences, with some favoring initial human brainstorming and others preferring direct AI interaction, indicating that usage patterns vary based on individual working styles.

At the same time, the study reveals notable challenges in usage behavior. Students often showed low prompting proficiency, failing to provide sufficient context, structure, or clarity in their queries. Moreover, many students were overly influenced by ChatGPT's confident responses, sometimes accepting incorrect or suboptimal suggestions without sufficient critique. This highlights a critical concern: while ChatGPT is perceived as useful, effective usage requires strong critical thinking and evaluation skills, as well as awareness of the risks of over-reliance.

\section{Implications}
The implications of this work suggest that conversational AI can be effectively integrated into software engineering education, but only when accompanied by appropriate guidance and pedagogical adjustments. The findings indicate that AI tools should not be treated as simple automation aids but rather as collaborative partners that require active human involvement. Teaching students how to construct effective prompts, critically evaluate AI outputs, and iteratively refine results becomes essential. This implies a shift in curriculum design toward emphasizing prompt engineering, critical thinking, and human–AI collaboration skills.

Another key implication is that different interaction modes with AI may suit different learners, and educational environments should accommodate this diversity rather than enforcing a single workflow. While AI can support productivity and idea generation, over-reliance on it without critical assessment can negatively impact learning outcomes. Therefore, educators should focus on developing students’ ability to question and validate AI-generated content. More broadly, the study highlights the need for a balanced approach where AI augments human reasoning without diminishing core problem-solving and analytical skills.

\section{Threats to Validity}
Threats to validity arise primarily from the study’s educational context and experimental design. The participants were undergraduate students in a single course at one institution, which may limit the generalizability of the findings to other populations such as professional developers or students with different backgrounds. Additionally, the study focuses specifically on requirements engineering tasks within a project-based course, and the results may not directly transfer to other software engineering activities. The controlled use of ChatGPT, including restricting students to a specific version and requiring submission of interaction logs, may also influence how students naturally interact with AI tools outside an academic setting.

Another important threat concerns the evaluation and measurement process. Although grading rubrics were carefully designed and cross-validated, assessing requirements quality and prompt effectiveness inherently involves some level of subjectivity. Variability in ChatGPT’s responses across different runs could also affect reproducibility, even though efforts were made to mitigate this through repeated experiments. Furthermore, students’ self-reported reflections and effort estimates may not always accurately capture their actual behavior, introducing potential bias in understanding interaction patterns and effort distribution.
